\section{Extensions: variance correction, a control variate, and faster sampling}
\label{sec:extensions}

Three weaknesses stand out in the results above: the low-rank sampler underprices at
every rank (§\ref{sec:errorrank}), plain Monte Carlo needs $\approx 130{,}000$
paths for a $1\%$ confidence half-width (§\ref{sec:convergence}), and the per-path
loops leave most of the machine idle: one core, one mat-vec at a time, and a random
number generator that dominates the FFT path (§\ref{sec:timing}).  This section adds
one remedy for each, and then uses the control variate to estimate Greeks.  All are
opt-in methods alongside the original pricers, measured by separate programs
(\texttt{./build/extensions}, \texttt{./build/performance}), so none of the results in
the preceding sections change.

\subsection{Variance-corrected low-rank sampler}
\label{sec:varcorr}

The low-rank sampler underprices because it discards the variance held by the small
eigenvalues of $C$.  The simplest repair puts that variance back at each time step as
independent noise:
\begin{equation}
  \label{eq:varcorr}
  W^H \approx L_k z + D^{1/2}\varepsilon, \qquad
  D = \mathrm{diag}\bigl(C - L_k L_k^\top\bigr), \quad
  z \sim \mathcal{N}(0, I_k),\ \varepsilon \sim \mathcal{N}(0, I_N).
\end{equation}
Every marginal variance $\mathrm{Var}(W^H(t_m))$ is then exact by construction, while the
off-diagonal covariances remain those of $C_k$.  The cost is $N$ extra Gaussian draws and
$N$ multiply-adds per path.

\begin{figure}[t]
  \centering
  \includegraphics[width=\linewidth]{variance_corrected_rank.png}
  \caption{Variance-corrected vs.\ plain low-rank sampler.
    \textbf{(a)} Price error against the reference ($\pm 2$ standard errors shaded).
    \textbf{(b)} Mean relative error of the increment (fGn) variances.
    \textbf{(c)} Mean relative error of the lag-1 increment covariances.
    Panels (b) and (c) are exact, computed from each sampler's covariance matrix.
    Controlled: $N = 500$, $H = 0.10$, $\nu = 0.52$, $\sigma_0 = 0.2$, $M = 100{,}000$
    paths per rank.
    Independent: rank $k$, sampler.
    Dependent: relative errors (\%).}
  \label{fig:varcorr}
\end{figure}

The correction removes the bias (Figure~\ref{fig:varcorr}(a)).  The plain sampler's
price error ranges from $-8.1\%$ at $k = 2$ to $-3.9\%$ at $k = 128$.  The corrected
sampler's error is within one standard error ($\approx 0.56\%$) at every rank, from
$-0.5\%$ to $+0.3\%$, for about 40\% more MC time at $k = 32$ (2.0~s to 2.8~s per
100{,}000 paths).

The key pitfall is that correct marginals do not mean correct paths.  The plain
low-rank paths are far too smooth: they miss 93--100\% of the increment variance
(Figure~\ref{fig:varcorr}(b)), which is where the roughness lives.  The correction's
noise is white in levels, so it overshoots and makes the increments too rough instead:
the increment-variance error is $119\%$ at $k = 2$, falls to $1.8\%$ at $k = 64$, and rises
again to $16\%$ at $k = 128$.  The lag-1 increment covariance improves from about $100\%$
wrong to $14$--$18\%$ wrong at $k \geq 64$.  The Asian price depends on the volatility
levels $\sigma_t$, not on their increments, which is why~\eqref{eq:varcorr} prices it
correctly.  A payoff that depends on path roughness itself, such as an option on
realized volatility, would need the residual's near-diagonal covariance as well, for
example through a banded Cholesky factor (§\ref{sec:futurework}).

\subsection{Conditional geometric control variate}
\label{sec:cv}

For constant volatility, the geometric-average Asian call has a closed-form price
\citep{kemna1990pricing} and is the classic control variate for the arithmetic one
\citep{glasserman2004monte}.  Under RFSV no such closed form exists, but one exists
\emph{conditionally} on the volatility path.  With $\rho = 0$ the price shocks are
independent of $\sigma$, so given the path, the log of the geometric average
$G = (\prod_{n} S_{t_n})^{1/N}$ is Gaussian with
\begin{equation}
  \label{eq:geomcond}
  m = \log S_0 + \Delta t \sum_{j=1}^{N} w_j\bigl(r - \tfrac{1}{2}\sigma_j^2\bigr),
  \qquad
  v = \Delta t \sum_{j=1}^{N} w_j^2 \sigma_j^2,
\end{equation}
where $w_j = (N - j + 1)/N$ is the share of the averaging window that the $j$-th shock
affects.  The
conditional expectation of $Y = (G - K)^+$ is then a Black--Scholes formula,
$\mathbb{E}[Y \mid \sigma] = e^{m + v/2}\,\Phi(d_1) - K\,\Phi(d_1 - \sqrt{v})$ with
$d_1 = (m - \log K + v)/\sqrt{v}$, computed per path in $O(N)$.

With $V$ the arithmetic payoff and $C = Y - \mathbb{E}[Y \mid \sigma]$, which has mean zero
exactly, the estimator
\begin{equation}
  \hat{p}_\mathrm{cv} = e^{-rT}\,\frac{1}{M}\sum_{m=1}^{M}\bigl(V^{(m)} - \beta\,C^{(m)}\bigr)
\end{equation}
is unbiased for any fixed $\beta$.  I estimate the variance-optimal coefficient
$\beta = \mathrm{Cov}(V, C)/\mathrm{Var}(C)$ on a separate pilot run of 10\% of the paths,
so the main estimate stays exactly unbiased; the variance is then reduced by the
factor $1/(1 - \mathrm{corr}(V, C)^2)$.  A single driver implements this for all three
samplers.

\begin{table}[t]
  \caption{Conditional geometric control variate ($M = 10{,}000$).  Time to reach a
    $0.1\%$ relative standard error is extrapolated from the measured standard error
    and per-path cost, including setup and, for the control variate, the pilot run.}
  \label{tab:cv}
  \centering
  \small
  \setlength{\tabcolsep}{5pt}
  \begin{tabular}{lrrrr}
    \toprule
    & \multicolumn{2}{c}{Variance reduction} & \multicolumn{2}{c}{Time to 0.1\% SE ($N = 1000$)} \\
    Sampler & $N = 252$ & $N = 1000$ & plain & control variate \\
    \midrule
    Cholesky             & 27.2 & 29.3 & 401~s & 16.9~s \\
    Circulant + FFT      & 28.8 & 28.1 & 234~s & 10.0~s \\
    rSVD ($k = 32$)      & 25.7 & 27.7 & 128~s & 5.7~s \\
    rSVD ($k = 32$) + diag & 26.6 & 25.8 & 179~s & 8.1~s \\
    \bottomrule
  \end{tabular}
\end{table}

\begin{figure}[t]
  \centering
  \includegraphics[width=\linewidth]{control_variate.png}
  \caption{Conditional geometric control variate.
    \textbf{(a)} Variance reduction $(\mathrm{SE}_\mathrm{plain}/\mathrm{SE}_\mathrm{cv})^2$
    for each sampler.
    \textbf{(b)} Time to a $0.1\%$ relative standard error at $N = 1000$, plain (light) and
    control variate (dark), on a log scale.
    Controlled: $H = 0.10$, $\nu = 0.52$, $\sigma_0 = 0.2$, $M = 10{,}000$.
    Independent: sampler, $N$, estimator.
    Dependent: variance reduction and time (s).}
  \label{fig:cv}
\end{figure}

Table~\ref{tab:cv} and Figure~\ref{fig:cv} show the results.  The variance falls by a
factor of 26--29 for every sampler at both grid sizes, and the time to a given accuracy
falls by about the same factor: at $N = 1000$, from 401~s to 17~s for Cholesky and from
234~s to 10~s for FFT.  For the purposes of this report, that is a larger saving than
any change of sampler: switching from Cholesky to FFT at $N = 1000$ gains $1.7\times$, while
the control variate gains about $23$--$24\times$ within either.

Two details limit the gain.  First, the reduction is set by
$\mathrm{corr}(V, C) = 0.981$, giving $1/(1 - 0.981^2) \approx 27$, not by the higher raw
correlation of the two payoffs (0.998).  Because $C$ has mean zero given the volatility
path, it cancels the noise from the price shocks but not the noise from the volatility
path itself, and that path-driven part is most of the roughly 4\% of variance that
remains: $\mathbb{E}[Y \mid \sigma]$ alone carries 2.7\% of $\mathrm{Var}(V)$.  A second
control on the volatility path, or quasi-Monte Carlo (§\ref{sec:qmc-future}), would be
needed to remove it.  Second, a smaller standard error does not fix bias.  The plain
rSVD sampler with the control variate converges $27\times$ faster to a price that is still
about 4\% too low; combined with the correction of §\ref{sec:varcorr}, it reaches an
accurate price in 8.1~s at $N = 1000$.

\paragraph{Limitations: }
The conditional closed form~\eqref{eq:geomcond} requires $\rho = 0$.  With spot-vol
correlation (§\ref{sec:rho-future}), conditioning on the volatility path no longer
leaves the price shocks independent, and the control variate would have to be
rebuilt, for instance by conditioning on the volatility's driving noise instead.

\subsection{Batched sampling, a faster generator, and threads}
\label{sec:perf}

The per-path loops of §\ref{sec:benchmarks} waste the hardware in three ways.  The
Cholesky mat-vec $Lz$ reads each element of $L$ once and uses it in a single
multiply-add, 2 flops per 8 bytes, so at large $N$ it runs at the speed of memory
rather than of the arithmetic units (§\ref{sec:memory}).  The Gaussian generator,
\texttt{std::mt19937} with \texttt{std::normal\_distribution}, costs 15~ns per draw and
dominates the FFT path.  And every path runs on one of eight cores.  I address each with
an opt-in \texttt{price\_batched()} in all three samplers, driven by one shared loop.

\paragraph{Design. }
Paths are generated in blocks of $B$.  Each sampler fills an $N \times B$ block of
log-volatility paths at once: Cholesky with one triangular matrix-matrix product $LZ$,
which reads $L$ once per block instead of once per path; the low-rank sampler with one
$N \times k$ by $k \times B$ product; the FFT sampler with one batched FFTW call over
$B/2$ transforms.  The fast generator combines xoshiro256++ \citep{blackman2021scrambled}
with a 256-layer ziggurat \citep{marsaglia2000ziggurat}, in which about 99\% of draws
cost one table comparison and one multiply; it draws in 3.8~ns, $4.0\times$ faster.
Blocks are also the unit of parallel work: threads take blocks from an atomic counter,
block $b$ draws from its own stream seeded by $(\text{seed}, b)$, and the block sums are
added in block order.  The price is therefore bitwise identical for any number of
threads, which the test suite checks.  Experiment design: controlled, $M = 10{,}000$
paths, seed, $k = 32$; independent, $N$, block size $B$, generator, and thread count;
dependent, Monte Carlo time (median of three repeats) and the time of each stage of a
path.

\begin{figure*}[t]
  \centering
  \includegraphics[width=\linewidth]{perf_breakdown.png}
  \caption{Time per path split into its three stages: drawing the normals, the
    transform (mat-vec, FFT or $L_k z$; normals replayed from a buffer), and the price
    path with its payoff.  For each sampler, the upper bar is the original per-path loop
    with \texttt{std::mt19937}, the lower bar the batched version ($B = 64$) with the
    ziggurat generator.  Each stage is timed in isolation, so the bars
    understate the end-to-end times of Figure~\ref{fig:perfbatched}.
    Controlled: $k = 32$, one thread.  Independent: $N$, sampler, configuration.
    Dependent: time per path ($\mu$s).}
  \label{fig:perfbreakdown}
\end{figure*}

\paragraph{Where the time goes. }
Figure~\ref{fig:perfbreakdown} splits each path into its stages.  The generator is a
fixed $4.0\times$ win wherever random numbers are drawn.  Batching helps only where the
transform is memory-bound: it cuts the Cholesky transform by $2.6\times$ at $N = 1000$
and $4.3\times$ at $N = 4000$ (from 1621 to 379~$\mu$s), and the low-rank product by
$2.6\times$.  At $N = 4000$ the batched product runs at $N^2/379\,\mu\mathrm{s}
\approx 42$~GFlop/s, roughly three quarters of the nominal double-precision peak of one
performance core (four 128-bit FMA units at 3.5~GHz), while the mat-vec streams $L$ at
$\approx 40$~GB/s.  The FFT transform does not change at all (5.4 against 5.3~$\mu$s at
$N = 1000$): each transform already works in cache, so a batch has nothing to reuse.
After both changes, the price path and its $2N$ exponentials, untouched here, are the
largest stage of a low-rank path and about a third of an FFT path.

\begin{figure}[t]
  \centering
  \includegraphics[width=\linewidth]{perf_batched.png}
  \caption{MC time per path vs $N$ for four configurations of each sampler, one thread:
    the original loop (dotted), unbatched with the ziggurat (dashed), and $B = 64$ with
    either generator (solid).  Color encodes the generator.  Bottom right: local scaling
    exponent of the Cholesky time between neighbouring $N$.
    Controlled: $M = 10{,}000$, $k = 32$, one thread.
    Independent: $N$, block size, generator.
    Dependent: MC time per path (median of 3 repeats; bars: interquartile range).}
  \label{fig:perfbatched}
\end{figure}

\paragraph{Speedups on one thread. }
Figure~\ref{fig:perfbatched} shows the end-to-end effect.  The FFT sampler gains
$2.0$--$2.2\times$ at every $N$, entirely from the generator.  The low-rank sampler gains
$1.7$--$1.9\times$, mostly from the generator.  The Cholesky gain grows with $N$, from
$2.2\times$ at $N = 64$ to $2.3\times$ at $N = 1000$ and $4.0\times$ at $N = 4000$
(19.3~s to 4.85~s), because batching removes the memory penalty: the local exponent
between $N = 2000$ and 4000 falls from 2.33 for the original loop to 1.74, below the
$O(N^2)$ flop count because the $O(N)$ stages still carry weight.  The block size has to
be large enough to pay off.  At $B = 2$ the triangular product is slower than the
mat-vec (127 against 87~$\mu$s per path at $N = 1000$), and the gain saturates around
$B = 64$.  Batching narrows the gap between the methods but does not reorder them: at
$N = 4000$ the batched FFT sampler is still $3.5\times$ faster than batched Cholesky (140
against 485~$\mu$s per path), down from $6.7\times$.

\begin{figure}[t]
  \centering
  \includegraphics[width=\linewidth]{perf_threads.png}
  \caption{Speedup over one thread vs thread count, unbatched (dashed; the FFT pairs two
    paths per transform) and $B = 64$ (solid), with the fast generator.  The shaded
    region uses the M2's efficiency cores.
    Controlled: $M = 10{,}000$ ($N = 500$) or $2{,}000$ ($N = 4000$), fast generator.
    Independent: threads, $N$, block size.
    Dependent: speedup of the MC time (median of 3 repeats).}
  \label{fig:perfthreads}
\end{figure}

\paragraph{Thread scaling as a test of the bandwidth argument. }
Figure~\ref{fig:perfthreads} gives the most direct evidence for the claim of
§\ref{sec:memory} that the large-$N$ Cholesky loop is memory-bound.  A compute-bound loop
should scale with cores, while a bandwidth-bound one should stall once the cores
saturate the memory bus.  At $N = 4000$ the unbatched Cholesky loop reaches only
$1.8\times$ (at three threads) and falls back to $1.5\times$ at eight, as extra threads
compete for the same bandwidth.  The batched version, which reads $L$ once per 64 paths,
reaches $3.1\times$ on the four performance cores and $4.1\times$ with all eight.  The FFT
and low-rank samplers scale similarly whether batched or not, to $3.3$--$3.8\times$ on
four cores at $N = 500$ and to $4.3$--$5.7\times$ on eight at either $N$.  Even at
$N = 500$, where $L$ fits in cache, the unbatched Cholesky loop scales worst
($2.7\times$ on four cores), since the performance cores share one L2.  The efficiency cores add
less than the performance cores, so no curve stays linear past four threads.

\paragraph{Limitations: }
The machine is a fanless laptop, so sustained runs can throttle: back-to-back
$N = 4000$ Cholesky runs slowed by $\approx 8\%$ after 40~s, and multithreaded runs heat
the chip faster.  The thread-scaling data come from a run with the machine otherwise
idle.  A repeat with a background process using one core gave the same qualitative
picture (unbatched Cholesky at $N = 4000$ peaking at $2.3\times$ and falling to
$1.6\times$), but noisier curves.  The stage times are measured in isolation with warm
caches, so their sum understates a full path by 20--40\%.  Batching changes which
paths a seed produces, so batched prices agree with the original ones only within Monte
Carlo error (checked in the test suite), not to the last digit.

\subsection{Greeks with common random numbers}
\label{sec:greeks}

With the control variate and common random numbers, sensitivities become cheap to
estimate precisely.  I compute two for the ATM Asian call at the model parameters:
vega with respect to the vol-of-vol $\nu$, and the sensitivity to the Hurst exponent
$H$ (\texttt{data/greeks.py}).  The point of comparing estimators is a consistency
check: two unrelated estimators of $\partial V/\partial\nu$ must agree.

The first is the central difference $[\hat V(\nu + h) - \hat V(\nu - h)]/(2h)$, with
both prices from the same random numbers.  Each path's payoff then moves continuously
with $\nu$, so the per-path difference is $O(h)$ and the variance stays bounded as
$h \to 0$; with independent random numbers it would grow like $h^{-2}$.  The second is
the pathwise estimator \citep{glasserman2004monte}, which needs no bump.  Since
$\partial\sigma_n/\partial\nu = \sigma_n W^H_{t_n}$, differentiating the discrete price
path gives
\begin{equation}
  \label{eq:pathwise}
  \frac{\partial V}{\partial \nu} = \mathbf{1}\{A > K\}\,\frac{1}{N}\sum_{n=1}^{N} S_n
  \sum_{m \leq n} \sigma_m W^H_{t_m}\bigl(\sqrt{\Delta t}\,Z_m - \sigma_m \Delta t\bigr),
\end{equation}
where $A$ is the arithmetic average.  The payoff is Lipschitz in $A$, so the derivative
and the expectation can be exchanged and~\eqref{eq:pathwise} is unbiased.  The control
variate differentiates too: $C = Y - \mathbb{E}[Y \mid \sigma]$ has mean zero for every
$\nu$, so $\partial C/\partial\nu$ does as well, and the closed
form~\eqref{eq:geomcond} differentiates through its mean and variance parameters.

\begin{figure}[t]
  \centering
  \includegraphics[width=\linewidth]{greeks.png}
  \caption{Greeks of the ATM Asian call.
    \textbf{(a)} $\partial V/\partial\nu$ by central difference with common random
    numbers and the control variate, against bump size $h$, and the pathwise
    estimator~\eqref{eq:pathwise} with and without its differentiated control.
    \textbf{(b)} $\partial V/\partial H$ by central difference.  Error bars are
    $\pm 1$ standard error from the per-path differences.
    Controlled: $H = 0.10$, $\nu = 0.52$, $\sigma_0 = 0.2$, $K = S_0 = 100$, $T = 1$,
    $r = 0$, $N = 252$, $M = 20{,}000$.
    Independent: bump size $h$, estimator.
    Dependent: $\partial V/\partial\nu$, $\partial V/\partial H$.}
  \label{fig:greeks}
\end{figure}

The two agree (Figure~\ref{fig:greeks}).  The pathwise vega with its control is
$2.900 \pm 0.035$; the central difference is $2.892 \pm 0.035$ at $h = 0.005$, $0.16$
standard errors away, and drifts up to $2.909$ at $h = 0.1$.  Since every bump size
reuses the same paths, that drift is the $O(h^2)$ curvature bias, measured more
precisely than the error bars suggest.  The differentiated control cuts the variance
of the pathwise estimator by a factor of 7, against 27 for prices.

The sensitivity to roughness is $\partial V/\partial H = -2.40 \pm 0.05$, flat for $h$
from $0.0025$ to $0.02$.  There is no simple pathwise form, because $H$ enters through
the square root of the circulant spectrum, but with common random numbers the bumped
paths move smoothly with $H$ and the finite difference is well behaved.  At the model
point, lowering $H$ by $0.01$ raises the price by $0.024$, $0.45\%$; raising $\nu$ by
$0.01$ raises it by $0.029$.  In these units the price is about as sensitive to
roughness as to vol-of-vol, consistent with the sensitivity surface of
Figure~\ref{fig:sensitivity}.
